\section{Mixed-event setup for TPC--BTOF matching}
\label{sec:me}

This section describes step 2: projecting every QA track to BTOF with the helix stored in the picoDst,
matching it at the cell level to the BTOF hits of the \emph{same} event (SE) and of \emph{other} events
(mixed events, ME), and validating the projection against the production matching.
Unless stated otherwise all numbers are from the same 10k-event test file as Secs.~\ref{sec:evtqa}--\ref{sec:tofqa}
(\meEvtSelTest{} selected events), produced by \fpath{StRoot/StTofMEMaker} and summarised by
\fpath{analysis_note/scripts/make_me.C} (\fpath{logs/me_numbers_test.txt}).

\subsection{Method}
\paragraph{BTOF geometry from data.} The picoDst does not contain the BTOF geometry, and the
official geometry (\fcmd{StBTofGeometry}) needs the full detector description plus database
alignment. Instead, the position of every module is fitted from production-matched tracks.
For these, \fcmd{btofHitPos()} is the production projection point in global coordinates and
($y_{\rm local}$, $z_{\rm local}$, cell) are its local coordinates, so for each module
\begin{equation}
\vec r_{\rm hit} = \vec C + y_{\rm local}\,\vec U + z_{\rm local}\,\vec V + (\mathrm{cell}-3.5)\,\vec W ,
\end{equation}
which is linear in the 12 parameters $(\vec C,\vec U,\vec V,\vec W)$. The maker accumulates the
normal-equation sums per module (\fcmd{hGeomSum}, additive over jobs) and
\fpath{macros/makeGeom.C} solves them (modules with $\ge20$ matched tracks).
Of 3840 modules, \geoNFitTest{} are fitted, with a mean residual of \geoFitRmsTest\,cm, i.e.\ the
picoDst position precision, so the linear model is exact. The fit gives $|\vec U|=\geoUmagTest$ (mean over modules), a cell pitch of
\geoPitchAbsTest\,cm (RMS \geoPitchRmsTest\,cm over modules), the \fcmd{mWidthPad} value of
\fcmd{StBTofMatchMaker}, and module-centre radii of \geoRMeanTest\,cm (RMS \geoRRmsTest\,cm;
Fig.~\ref{fig:me_test_geo_r}). The remaining \geoNRotTest{} modules (dead in this run or with too few tracks) are
copied from the same module of the nearest fitted tray in the same half (east/west), rotated about $z$
by the tray-to-tray $\phi$ difference. Applying the same rotation to fitted modules (closure test) reproduces
their centres to \geoClosureMeanTest\,cm on average (RMS \geoClosureRmsTest\,cm, at most
\geoClosureMaxTest\,cm; Fig.~\ref{fig:me_test_geo_closure}), small compared with the 3.45\,cm pitch.
Only \meFracRotGeomTest\% of the projected tracks land in a rotated module.

\choice{Data-driven module geometry instead of \fcmd{StBTofGeometry}: it includes the production
alignment automatically, has no database dependence, and is validated to the 0.01\,cm level by the
fit residuals. Geometry file: \fpath{geom/btof_geom_test.txt} (from the test run 20093036).}

\paragraph{Projection.} For each QA track (Sec.~\ref{sec:trkcuts}) the global helix is built with
\fcmd{StPicoTrack::helix(B)} (origin = DCA point, global momentum, charge, $B$ from the event).
It is intersected with a cylinder of radius 213.5\,cm (the mean module radius) to select candidate
modules ($|\Delta\phi|<0.09$ to the tray centre, $|\Delta z|<12$\,cm to the module centre); then the exact
intersection with each candidate module plane gives the local coordinates, the cell (pitch 3.45\,cm) and
$z_{\rm local}$. The projection is accepted if it falls inside one of the six cells and $|z_{\rm local}|<3.05$\,cm
(the production window). If several modules qualify, the one with smallest $|z_{\rm local}|$ is kept.

\paragraph{Matching definitions.} Hits are stored only as cell IDs, so all matches are at the cell level:
\begin{itemize}
\item \textbf{cell}: a BTOF hit in the projected cell;
\item \textbf{cell$\pm$1 (production-like)}: a hit in the projected cell or in an adjacent cell of the same module.
  This is the production rule: \fcmd{StBTofMatchMaker} (SL24y, ``match find neighbours'') accepts a hit in the
  projected cell or its $\pm1$ neighbour (checked in the source).
\end{itemize}
$P_{\rm SE}$ is the fraction of projected tracks matched to a hit of the same event.
$P_{\rm ME}$ is the fraction of (track, partner-event) trials in which the projected cell (or cell$\pm$1) has a hit
in the partner event. $f_{\rm rand}=P_{\rm ME}/P_{\rm SE}$ estimates the fraction of SE matches that are random.
To first order the true matching efficiency is $\varepsilon_{\rm true}=(P_{\rm SE}-P_{\rm ME})/(1-P_{\rm ME})$.

\paragraph{Mixing.} Partner events come from a FIFO buffer of the previous 5 selected events with the
same run number, the same $V_z$ class (14 classes of 10\,cm) and the same refMult class (boundaries 0, 10, 25,
50, 80, 120, 170, 230, 300). Buffers are cleared at each run change, so dead areas are identical in mixed pairs.
\meFracFullMixTest\% of the events have the full 5 partners.

\choice{Mixing classes in refMult (not number of BTOF hits) and $V_z$, depth 5, same run only.
refMult is independent of BTOF, and the hit occupancy follows it closely (Fig.~\ref{fig:evt_refmult_vs_tofmult});
$P_{\rm ME}$ is then studied directly versus the hit multiplicity of the partner event.}

\subsection{Validation of the projection}
Of the \meNTrkQATest{} QA tracks, \meFracProjOkTest\% have an accepted projection; the rest have no crossing
with the TOF radius (\meFracNoRTest\%), fall between modules (\meFracNoModTest\%), outside the six cells
(\meFracOutYTest\%) or outside $|z_{\rm local}|<3.05$\,cm (\meFracOutZTest\%) (Fig.~\ref{fig:me_test_projstatus}).

For the \meNProdMatchTest{} production-matched QA tracks, \meFracPMProjTest\% have an accepted projection,
\meFracPMModTest\% are projected to the same module as production and \meFracPMCellTest\% to the same cell
(Fig.~\ref{fig:me_test_prodcmp}). Above $p=1.5$\,\GeVc{} the values are \meFracPMModHiTest\% (module) and
\meFracPMCellHiTest\% (cell); for $0.2<p<0.4$\,\GeVc{} the same-cell fraction is only \meFracPMCellLoTest\%.
Different cells in the same module are expected: \meProdYlOutTest\% of the production matches have
$|y_{\rm local}|>1.75$\,cm, i.e.\ production matched them to a hit in the neighbouring cell.
When the cells agree, the local-coordinate differences (Figs.~\ref{fig:me_test_dy},~\ref{fig:me_test_dz}) have
RMS \meDyRmsHiTest\,cm ($y$) and \meDzRmsHiTest\,cm ($z$) for $p>1.5$\,\GeVc, growing to
\meDyRmsLoTest\,cm and \meDzRmsLoTest\,cm for $p<0.45$\,\GeVc.

\caveat{The picoDst helix starts at the DCA to the primary vertex, while production projects the outer
TPC helix. Energy loss and multiple scattering between the TPC and BTOF are therefore not in this projection,
which is why the agreement degrades at low momentum and $P_{\rm SE}$ is below the production match rate
for $p_T<0.3$\,\GeVc{} (Table~\ref{tab:me}). $P_{\rm ME}$ is not affected (Sec.~\ref{sec:meres}).}

\begin{figure}[H]
\centering
\figpanel{0.32}{me_test_geo_r}{Radius of the fitted module centres.}
\hfill
\figpanel{0.32}{me_test_geo_closure}{Closure test of the rotated geometry: distance between rotated and
fitted module centres, log scale.}
\hfill
\figpanel{0.32}{me_test_projstatus}{Projection status of QA tracks (fractions).}
\\[1ex]
\figpanel{0.32}{me_test_prodcmp}{For production-matched tracks: fraction with an accepted projection (grey),
projected to the production module (blue), and to the production cell (red), vs $p$.}
\hfill
\figpanel{0.32}{me_test_dy}{$y_{\rm local}^{\rm proj}-y_{\rm local}^{\rm prod}$ vs $p$ when the cells agree, log colour scale.}
\hfill
\figpanel{0.32}{me_test_dz}{$z_{\rm local}^{\rm proj}-z_{\rm local}^{\rm prod}$ vs $p$ when the cells agree, log colour scale.}
\caption{Geometry and projection validation (test sample).}
\end{figure}

\subsection{Same-event vs mixed-event matching (test sample)}
\label{sec:meres}
\meNProjTest{} QA tracks are projected and tested against \meNTrialsTest{} partner events in total.
Integrated over $p_T$:
\begin{center}
\begin{tabular}{lccc}
\toprule
 & $P_{\rm SE}$ & $P_{\rm ME}$ (\%) & $f_{\rm rand}=P_{\rm ME}/P_{\rm SE}$ (\%) \\
\midrule
cell & \mePSETest & \mePMEpctTest & \meFakeFracTest \\
cell$\pm$1 (production-like) & \mePSENbTest & \mePMENbpctTest & \meFakeFracNbTest \\
\bottomrule
\end{tabular}
\end{center}
The production match rate for the same projected tracks is \mePProdTest.
The statistical uncertainty on $P_{\rm ME}$ is $\le0.01$\% (absolute).

\begin{table}[h]
\centering
\caption{Matching probabilities vs $p_T$ for projected QA tracks (generated table \fpath{tables/me_eff_test.tex}).
$P_{\rm prod}$: production match with $\beta>0$. Statistical errors on $P_{\rm ME}$; those on $P_{\rm SE}$ and $P_{\rm prod}$ are below 0.003.}
\label{tab:me}
\small
\input{tables/me_eff_test.tex}
\end{table}

Main observations:
\begin{itemize}
\item \textbf{$P_{\rm ME}$ is the hit occupancy.} $P_{\rm ME}$ grows linearly with the number of hits $N_{\rm hits}$ in the partner
  event (Fig.~\ref{fig:me_test_prob_nh}). A fit $P_{\rm ME}=a\,N_{\rm hits}$ gives $a=\meSlopePerHitTest\times10^{-5}$,
  which is \meSlopeRatioTest{} times the uniform-occupancy expectation $1/23040$. Averaged over trials,
  $P_{\rm ME}/(\langle N_{\rm hits}\rangle/23040)=\mePMEoverOccTest$.
\item \textbf{$P_{\rm ME}$ is nearly flat} (2.3--2.5\% for the cell definition) in $p_T$, $\eta$ and in the local coordinates
  (Figs.~\ref{fig:me_test_prob_pt}, \ref{fig:me_test_prob_eta}, \ref{fig:me_test_yl}, \ref{fig:me_test_zl}),
  as expected for hits unrelated to the track; it rises with refMult (Fig.~\ref{fig:me_test_prob_ref}).
  Because random hits do not depend on the track, $P_{\rm ME}$ does not depend on the projection precision.
\item \textbf{Random-match fraction.} For the production-like definition $f_{\rm rand}$ is \meFakeMidNbTest\% for
  $0.5<p_T<2$\,\GeVc{} and \meFakeLowNbTest\% for $0.15<p_T<0.3$\,\GeVc{} (Fig.~\ref{fig:me_test_fake_pt}); the
  increase at low $p_T$ comes from the lower $P_{\rm SE}$ there.
\item \textbf{Same-event shape.} $P_{\rm SE}$ falls towards the cell edges in $y_{\rm local}$
  (\mePSEylCoreTest{} for $|y|<1$\,cm, \mePSEylEdgeTest{} for $1.3<|y|<1.6$\,cm) and at $|z_{\rm local}|\gtrsim2.8$\,cm;
  there the true hit is often in the neighbouring cell or module.
\item First-order true matching efficiency: $\varepsilon_{\rm true}=\meTrueEffNbTest$ (cell$\pm$1) and
  \meTrueEffTest{} (cell).
\end{itemize}

\caveat{$f_{\rm rand}$ treats every projected track as equally exposed to random hits. Tracks whose own hit
is present can still pick up a random hit in a neighbouring cell, which production would resolve differently.
Production also resolves hits claimed by several tracks (flag 2) and requires a valid time, while $P_{\rm SE}$
here does not; this is why $P_{\rm SE}^{\pm1}$ exceeds $P_{\rm prod}$ above 0.3\,\GeVc.}

\begin{figure}[H]
\centering
\figpanel{0.48}{me_test_prob_pt}{Matching probabilities vs $p_T$ (log scale): production (open black),
SE cell$\pm$1 (open blue), SE cell (blue), ME cell$\pm$1 (open magenta), ME cell (red).}
\hfill
\figpanel{0.48}{me_test_fake_pt}{$f_{\rm rand}=P_{\rm ME}/P_{\rm SE}$ vs $p_T$ (\%) for cell$\pm$1 (open magenta) and cell (red).}
\\[1ex]
\figpanel{0.48}{me_test_prob_nh}{$P_{\rm ME}$ (cell) vs number of BTOF hits in the partner event (red);
dashed blue: linear fit through the origin; dotted grey: uniform occupancy $N_{\rm hits}/23040$.}
\hfill
\figpanel{0.48}{me_test_prob_ref}{$P_{\rm SE}$ (blue) and $P_{\rm ME}$ (red), cell definition, vs refMult of the track's event.}
\\[1ex]
\figpanel{0.32}{me_test_prob_eta}{$P_{\rm SE}$ and $P_{\rm ME}$ (cell) vs $\eta$.}
\hfill
\figpanel{0.32}{me_test_yl}{$P_{\rm SE}$ and $P_{\rm ME}$ (cell) vs projected $y_{\rm local}$.}
\hfill
\figpanel{0.32}{me_test_zl}{$P_{\rm SE}$ and $P_{\rm ME}$ (cell) vs projected $z_{\rm local}$.}
\caption{Same-event (blue) and mixed-event (red) matching, test sample. Errors: 68\% binomial intervals.}
\end{figure}
